%! Author = sriadityavedantam
%! Date = 3/30/26

\section{Congruences and Lagrange's}

\begin{definition}[Left Congruence]
    Let $H\leq G$.
    Define a relation on $G$ by
    \[
        a\sim b,\qquad a\overset{\ell}{\sim} b\mod H
    \]
    if
    \[
        \vs{a}b\in H.
    \]
    Equivalently,
    \[
        \vs{b}a\in H.
    \]
    This is called left congruence modulo $H$.
\end{definition}

\begin{theorem}{
    Left congruence is an equivalence relation.
}
    \textbf{1)}
    It is reflexive.
    If $a\in G$, then
    \[
        \vs{a}a=e\in H.
    \]
    \textbf{2)}
    Let $a,b\in G$.
    Then
    \[
        a\sim b \iff \vs{a}b\in H.
    \]
    Since $H$ is closed under inverses,
    \[
        \vs{a}b\in H \iff (\vs{a}b)^{-1}\in H \iff \vs{b}a\in H.
    \]
    Thus
    \[
        a\sim b \iff b\sim a.
    \]
    Hence $\sim$ is symmetric.

    \textbf{3)}
    Suppose $a,b,c\in G$ and $a\sim b$, $b\sim c$.
    Then
    \[
        \vs{a}b\in H
        \quad\text{and}\quad
        \vs{b}c\in H.
    \]
    So
    \[
        (\vs{a}b)(\vs{b}c)=\vs{a}c\in H.
    \]
    Thus
    \[
        a\sim c.
    \]
    Hence $\sim$ is transitive.
\end{theorem}

What are the left equivalence classes?

\begin{definition}[Left Congruence Classes]
    The left congruence class of $a$ is
    \[
        \{b\in G:b\equiv a\mod H\}
        =
        \{b\in G:\exists h\in H,\ b=ah\}
        =
        aH
        =
        \{ah:h\in H\}.
    \]
    We call $aH$ the left coset containing $a$.
    These partition $G$.
    That is, $G$ is the union of its left cosets, and any two left cosets are either identical or disjoint.
\end{definition}

Suppose
\[
    G=S_3
    \quad\text{and}\quad
    H=\gen{(12)}=\{e,(12)\}.
\]
Then
\begin{gather*}
    (13)H=\{(13),(132)\},\\
    (23)H=\{(23),(123)\}.\\
\end{gather*}
Thus
\[
    S_3=\{e,(12)\}\cup\{(13),(132)\}\cup\{(23),(123)\}.
\]
\begin{proposition}{
    All left cosets of $H$ have $\#H$ elements.
}
    Let $a\in G$.
    Then
    \[
        aH=\{ah:h\in H\}
    \]
    is a left coset of $H$ in $G$.
    Define a map
    \[
        \phi:H\to aH
    \]
    by
    \[
        \phi(h)=ah.
    \]
    This map is injective, since
    \[
        ah_1=ah_2 \implies h_1=h_2
    \]
    by cancellation in $G$.
    It is also surjective by definition of $aH$.
    So $\phi$ is a bijection.
    Hence
    \[
        \#H=\#aH.
    \]
    Note that $aH$ need not be a group.
    This is only a set bijection.
\end{proposition}

\begin{definition}[Index of $H$ in $G$]
    The number of distinct left cosets of $H$ in $G$ is called the index of $H$ in $G$, denoted by
    \[
        [G:H].
    \]
    If $G$ is finite, then
    \[
        [G:H]=\frac{\ord(G)}{\ord(H)}.
    \]
\end{definition}

\begin{theorem}[Lagrange's Theorem]{
    Suppose $G$ is a finite group.
    If $H\leq G$, then
    \[
        \ord(H)\mid \ord(G).
    \]
}
    Suppose
    \[
        G=a_{1}H\cup\dots\cup a_{k}H,
    \]
    where the $a_{i}H$ are the distinct left cosets of $H$ in $G$.
    Then
    \[
        k=[G:H].
    \]
    Each coset has $\ord(H)$ elements.
    So
    \[
        \ord(G)=k\ord(H).
    \]
    Therefore
    \[
        \ord(H)\mid \ord(G).
    \]
\end{theorem}

\begin{corollary}[Corollary of Lagrange's Theorem]{
    Suppose $G$ is a finite group of prime order.
    If $H\leq G$, then
    \[
        H=\{e\}
        \quad\text{or}\quad
        H=G.
    \]
}
    Since $\ord(H)\mid \ord(G)$ and $\ord(G)=p$ is prime, we must have
    \[
        \ord(H)=1
        \quad\text{or}\quad
        \ord(H)=p.
    \]
    Hence
    \[
        H=\{e\}
        \quad\text{or}\quad
        H=G.
    \]
\end{corollary}

\begin{corollary}[Corollary of the Previous Corollary]{
    Suppose $G$ is a finite group of prime order.
    Then $G$ is cyclic.
    In fact,
    \[
        G\cong (\zp,+).
    \]
}
    Let $1\neq a\in G$.
    Let
    \[
        H=\gen{a}.
    \]
    Then $H$ is a subgroup of $G$ and is not trivial.
    By the previous corollary, we must have
    \[
        H=G.
    \]
    So $G$ is cyclic.
    Since $\ord(G)=p$, we get
    \[
        G\cong (\zp,+).
    \]
\end{corollary}

\begin{corollary}[Corollary of Fermat's Little Theorem]{
    Let $p$ be prime, $a\in\z$, and $p\nmid a$.
    Then
    \[
        a^{p-1}\equiv 1\mod p.
    \]
}
    Let
    \[
        G=\zp^*.
    \]
    Then $[a]\in \zp^*$.
    By Lagrange's Theorem,
    \[
        \ord(\gen{[a]})\mid (p-1).
    \]
    Hence
    \[
        [a]^{p-1}=[1]
    \]
    in $\zp^*$.
    Therefore
    \[
        a^{p-1}\equiv 1\mod p.
    \]
\end{corollary}

\begin{corollary}[Corollary 2 of Fermat's Little Theorem]{
    For all $a\in\z$,
    \[
        a^p\equiv a\mod p.
    \]
    Thus
    \[
        \frac{a^p-a}{p}\in\z.
    \]
}
    If $p\mid a$, then
    \[
        a\equiv 0\mod p,
    \]
    so
    \[
        a^p\equiv 0\equiv a\mod p.
    \]
    If $p\nmid a$, then by Fermat's Little Theorem,
    \[
        a^{p-1}\equiv 1\mod p.
    \]
    Multiplying by $a$ gives
    \[
        a^p\equiv a\mod p.
    \]
    Thus in all cases,
    \[
        p\mid (a^p-a),
    \]
    so
    \[
        \frac{a^p-a}{p}\in\z.
    \]
\end{corollary}

\begin{example}
    Suppose $\ord(G)=4$.
    Let $1\neq a\in G$.
    Since
    \[
        \ord(\gen{a})\mid 4,
    \]
    we have
    \[
        \ord(\gen{a})=2
        \quad\text{or}\quad
        \ord(\gen{a})=4.
    \]
    If $\ord(a)=4$, then $G$ is cyclic and
    \[
        G\cong \z_4.
    \]
    Suppose instead that every nonidentity element has order $2$.
    Let $1\neq a,b\in G$ with $a\neq b$.
    Then
    \[
        G=\{1,a,b,ab\}.
    \]
    In this case,
    \[
        G\cong C_2\times C_2,
    \]
    the Klein-$4$ group.
\end{example}